\begin{abstract}
For patients managing several medicines, verification of loose pills involves recognising newly enrolled medicines while rejecting pills outside a personal medication list. PharmaScan formulates this task as retrieval from a personal gallery, with explicit outcomes for an accepted match, an unknown pill and a request for another photograph. A hybrid visual encoder is trained with an angular margin objective to produce a 256-dimensional embedding, through which each medicine is enrolled from three photographs without network retraining. Rejection and ambiguity thresholds are selected using separate validation medicines, while optical character recognition with Florence-2 supplies complementary imprint evidence. On 200 held-out appearance groups from the Computational Photography Project for Pill Identification (C3PI) collection, simulated galleries of ten medicines yield 76.0\% correct naming and 0.04\% wrong naming among enrolled queries, with 94.8\% rejection of medicines outside the gallery. Imprint re-ranking increases identification accuracy among 4,936 catalogue products from 71.6\% to 80.2\%. Transfer to MEDISEG reveals a substantial acquisition gap, with 70.4\% of enrolled queries rejected as unknown. These findings support calibrated verification under familiar imaging conditions and identify acquisition shift and imprint visibility as constraints on broader use.
\end{abstract}
